\documentclass[10pt,a4paper]{amsart}
\usepackage{amsmath,amssymb,amsthm,graphicx}
% Retained minimal preamble for the Stacks Project chapter more-algebra.tex (More on Algebra).
% The chapter file pinned-original-source.tex carries only \input{preamble} on its first line, so no
% preamble text lives inside the chapter itself.  The Stacks Project uses one single shared preamble for
% every chapter; the snapshot stages it identically as sources/stacks-chow/preamble.tex and
% sources/stacks-etale-cohomology/preamble.tex.  The block below is copied mechanically, line for line,
% from sources/stacks-chow/preamble.tex lines 14-18 (the Xy-pic setup) and 154-217 (the native theorem
% declarations together with the mathematical macros).  Nothing was added, removed or redefined; the
% wrapper class and the amsmath/amssymb/amsthm/graphicx base packages are supplied by the assembler.
% Chapter: more-algebra.tex -- More on Algebra.
% For commutative diagrams we use Xy-pic
\usepackage[all]{xy}

% We use 2cell for 2-commutative diagrams.
\xyoption{2cell}

% Theorem environments.
%
\theoremstyle{plain}
\newtheorem{theorem}[subsection]{Theorem}
\newtheorem{proposition}[subsection]{Proposition}
\newtheorem{lemma}[subsection]{Lemma}

\theoremstyle{definition}
\newtheorem{definition}[subsection]{Definition}
\newtheorem{example}[subsection]{Example}
\newtheorem{exercise}[subsection]{Exercise}
\newtheorem{situation}[subsection]{Situation}

\theoremstyle{remark}
\newtheorem{remark}[subsection]{Remark}
\newtheorem{remarks}[subsection]{Remarks}

\numberwithin{equation}{subsection}

% Macros
%
\def\lim{\mathop{\mathrm{lim}}\nolimits}
\def\colim{\mathop{\mathrm{colim}}\nolimits}
\def\Spec{\mathop{\mathrm{Spec}}}
\def\Hom{\mathop{\mathrm{Hom}}\nolimits}
\def\Ext{\mathop{\mathrm{Ext}}\nolimits}
\def\SheafHom{\mathop{\mathcal{H}\!\mathit{om}}\nolimits}
\def\SheafExt{\mathop{\mathcal{E}\!\mathit{xt}}\nolimits}
\def\Sch{\mathit{Sch}}
\def\Mor{\mathop{\mathrm{Mor}}\nolimits}
\def\Ob{\mathop{\mathrm{Ob}}\nolimits}
\def\Sh{\mathop{\mathit{Sh}}\nolimits}
\def\NL{\mathop{N\!L}\nolimits}
\def\CH{\mathop{\mathrm{CH}}\nolimits}
\def\proetale{{pro\text{-}\acute{e}tale}}
\def\etale{{\acute{e}tale}}
\def\QCoh{\mathit{QCoh}}
\def\Ker{\mathop{\mathrm{Ker}}}
\def\Im{\mathop{\mathrm{Im}}}
\def\Coker{\mathop{\mathrm{Coker}}}
\def\Coim{\mathop{\mathrm{Coim}}}

% Boxtimes
%
\DeclareMathSymbol{\boxtimes}{\mathbin}{AMSa}{"02}

%
% Macros for moduli stacks/spaces
%
\def\QCohstack{\mathcal{QC}\!\mathit{oh}}
\def\Cohstack{\mathcal{C}\!\mathit{oh}}
\def\Spacesstack{\mathcal{S}\!\mathit{paces}}
\def\Quotfunctor{\mathrm{Quot}}
\def\Hilbfunctor{\mathrm{Hilb}}
\def\Curvesstack{\mathcal{C}\!\mathit{urves}}
\def\Polarizedstack{\mathcal{P}\!\mathit{olarized}}
\def\Complexesstack{\mathcal{C}\!\mathit{omplexes}}
% \Pic is the operator that assigns to X its picard group, usage \Pic(X)
% \Picardstack_{X/B} denotes the Picard stack of X over B
% \Picardfunctor_{X/B} denotes the Picard functor of X over B
\def\Pic{\mathop{\mathrm{Pic}}\nolimits}
\def\Picardstack{\mathcal{P}\!\mathit{ic}}
\def\Picardfunctor{\mathrm{Pic}}
\def\Deformationcategory{\mathcal{D}\!\mathit{ef}}

% --- appended 2026-10-12 (Chief): the project own definitions of the reference/slogan/history
% display environments, copied mechanically from the shared Stacks preamble lines 8-12
% (\usepackage{verbatim} plus the three \newenvironment lines, which the chapter preamble keeps
% verbatim from the project preamble). Nothing added, removed or redefined: these environments
% carry no mathematical content and the project itself typesets them as comments. Without these
% three lines a counted statement whose own text uses them fails with Environment slogan undefined.
\usepackage{verbatim}
\newenvironment{reference}{\comment}{\endcomment}
\newenvironment{slogan}{\comment}{\endcomment}
\newenvironment{history}{\comment}{\endcomment}
\usepackage{cleveref}
\crefname{originallemma}{original lemma}{original lemmas}
\Crefname{originallemma}{Original lemma}{Original lemmas}
\makeatletter
\newlabel{algebra-proposition-dimension-zero-ring}{{00KJ}{source}{}{sourcepointer-1}{}}
\newlabel{algebra-proposition-dimension-zero-ring@cref}{{[originallemma][0][]00KJ}{[1][0][]source}{}{}{}}
\makeatother
\crefname{originallemma}{original lemma}{original lemmas}
\Crefname{originallemma}{Original lemma}{Original lemmas}
\makeatletter
\newlabel{lemma-generalized-valuation-ring}{{37638}{source}{}{sourcepointer-2}{}}
\newlabel{lemma-generalized-valuation-ring@cref}{{[originallemma][0][]37638}{[1][0][]source}{}{}{}}
\makeatother
\crefname{originallemma}{original lemma}{original lemmas}
\Crefname{originallemma}{Original lemma}{Original lemmas}
\makeatletter
\newlabel{lemma-generalized-valuation-ring-modules}{{37675}{source}{}{sourcepointer-3}{}}
\newlabel{lemma-generalized-valuation-ring-modules@cref}{{[originallemma][0][]37675}{[1][0][]source}{}{}{}}
\makeatother
\crefname{originallemma}{original lemma}{original lemmas}
\Crefname{originallemma}{Original lemma}{Original lemmas}
\makeatletter
\newlabel{lemma-split-off-free-part}{{37847}{source}{}{sourcepointer-4}{}}
\newlabel{lemma-split-off-free-part@cref}{{[originallemma][0][]37847}{[1][0][]source}{}{}{}}
\makeatother
\title{Every finite module over a principal ideal domain is a finite free module plus a sum of quotients by nonzero elements}
\author{The Stacks Project Authors}
\date{stacks-project:more-algebra.tex:local-snapshot-20261012; GNU FDL 1.2 or later}
\begin{document}
\maketitle
\setcounter{subsection}{1}
\noindent\textit{Editorial note.} All delivered mathematical text is the author's own, copied from the pinned Stacks Project chapter file more-algebra.tex; no statement or proof step was rewritten, reordered or removed.\par\medskip
\noindent\textit{Editorial note.} The chapter file pinned-original-source.tex carries only \textbackslash{}input\{preamble\} on its first line, so no preamble text lives inside the chapter. The retained file sources/stacks-more-algebra/retained-preamble.tex is a mechanical, line-for-line copy of the project's single shared preamble as staged in the snapshot; nothing was added, removed or redefined and no package was added.\par\medskip
\noindent\textit{Editorial note.} The source URL is the source-level Stacks Project URL, as required when more than one card is produced from one source; the per-claim tag is carried by claim\_key.\par\medskip
\noindent\textit{Editorial note.} The author writes '(small detail omitted)' twice inside the counted proof span and writes 'Since \$R/fR\$ is Noetherian of dimension \$0\$ (small detail omitted) we see that \$R/fR = \textbackslash{}prod R\_j\$ is a finite product of Artinian local rings \$R\_i\$', using j and i inconsistently; both omissions and the inconsistent index are reproduced unchanged. The four references inside the counted proof span are recorded as non-load-bearing navigation pointers.\par\medskip
\noindent\textit{Allowed standard prerequisites (editorial).}\begin{itemize}
\item The splitting of a finite module over a Bezout domain into a free part and a torsion part. invoked at proof lines 37897-37898. Reduces the statement to the torsion case.
\item A zero-dimensional Noetherian ring is a finite product of Artinian local rings. invoked at proof lines 37901-37904. Produces the factors R\_j over which the module decomposes.
\item The structure of finite modules over a generalized valuation ring. invoked at proof lines 37910-37911. Gives each M\_j as a direct sum of cyclic quotients over R\_j.
\item Lifting cyclic quotients: choosing g\_\{ji\} with (g\_\{ji\}) = (f\_j, f\_\{ji\}). used at proof lines 37912-37913. Is the step that turns the decomposition over R/fR into one over R.
\end{itemize}\medskip
\ifdefined\hypertarget\hypertarget{sourcepointer-1}{}\fi
\noindent\textit{Source reference.} Original lemma 00KJ: focused-seeds/AGP/sources/stacks-snapshot-20261011/tags.txt, source lines 753--753. Cross-chapter navigation pointer into the Stacks Project; the standard fact it supplies is applied inside the counted proof span and the argument resting on it stays there.\par\medskip
\ifdefined\hypertarget\hypertarget{sourcepointer-2}{}\fi
\noindent\textit{Source reference.} Original lemma 37638: sources/stacks-more-algebra/pinned-original-source.tex, source lines 37638--37638. Same-chapter navigation pointer; the statement it names is standard for the section and the argument that consumes it stays inside the counted proof span.\par\medskip
\ifdefined\hypertarget\hypertarget{sourcepointer-3}{}\fi
\noindent\textit{Source reference.} Original lemma 37675: sources/stacks-more-algebra/pinned-original-source.tex, source lines 37675--37675. Same-chapter navigation pointer; the statement it names is standard for the section and the argument that consumes it stays inside the counted proof span.\par\medskip
\ifdefined\hypertarget\hypertarget{sourcepointer-4}{}\fi
\noindent\textit{Source reference.} Original lemma 37847: sources/stacks-more-algebra/pinned-original-source.tex, source lines 37847--37847. Same-chapter navigation pointer; the statement it names is standard for the section and the argument that consumes it stays inside the counted proof span.\par\medskip
\begin{lemma}
\label{lemma-modules-PID}
Let $R$ be a PID. Every finite $R$-module $M$ is of isomorphic
to a module of the form
$$
R^{\oplus r} \oplus \bigoplus\nolimits_{i = 1, \ldots, n} R/f_iR
$$
for some $r, n \geq 0$ and $f_1, \ldots, f_n \in R$ nonzero.
\end{lemma}

\begin{proof}
A PID is a Noetherian B\'ezout ring. By Lemma \ref{lemma-split-off-free-part}
it suffices to prove the result if $M$ is torsion. Since $M$ is finite, this
means that the annihilator of $M$ is nonzero. Say $fM = 0$ for some
$f \in R$ nonzero. Then we can think of $M$ as a module over $R/fR$.
Since $R/fR$ is Noetherian of dimension $0$ (small detail omitted)
we see that $R/fR = \prod R_j$ is a finite product of Artinian
local rings $R_i$
(Algebra, Proposition \ref{algebra-proposition-dimension-zero-ring}).
Each $R_i$, being a local ring and a quotient of a PID, is a generalized
valuation ring in the sense of
Lemma \ref{lemma-generalized-valuation-ring} (small detail omitted).
Write $M = \prod M_j$ with $M_j = e_j M$ where $e_j \in R/fR$ is
the idempotent corresponding to the factor $R_j$.
By Lemma \ref{lemma-generalized-valuation-ring-modules}
we see that $M_j = \bigoplus_{i = 1, \ldots, n_j} R_j/\overline{f}_{ji}R_j$
for some $\overline{f}_{ji} \in R_j$. Choose lifts $f_{ji} \in R$
and choose $g_{ji} \in R$ with $(g_{ji}) = (f_j, f_{ji})$.
Then we conclude that
$$
M \cong \bigoplus R/g_{ji}R
$$
as an $R$-module which finishes the proof.
\end{proof}

\end{document}
